% ============================================================
%  Tactical Formation — Draft Paper
%  Compiler : XeLaTeX  (Overleaf: Menu > Compiler > XeLaTeX)
% ============================================================
\documentclass[12pt, a4paper]{article}

% ── Encoding & Language ──────────────────────────────────────
% Hỗ trợ cả pdfLaTeX và XeLaTeX:
%   - pdfLaTeX: dùng inputenc + T5 + babel (mặc định Overleaf)
%   - XeLaTeX : bỏ comment 3 dòng dưới, comment 3 dòng trên
\usepackage[utf8]{inputenc}
\usepackage[T5]{fontenc}
\usepackage[vietnamese]{babel}
% \usepackage{fontspec}        % Bỏ comment nếu dùng XeLaTeX
% \usepackage{polyglossia}     % Bỏ comment nếu dùng XeLaTeX
% \setmainlanguage{vietnamese} % Bỏ comment nếu dùng XeLaTeX

% ── Page Layout ──────────────────────────────────────────────
\usepackage[top=2.5cm, bottom=2.5cm, left=2.5cm, right=2.5cm]{geometry}

% ── Math ─────────────────────────────────────────────────────
\usepackage{amsmath, amssymb}

% ── Tables ───────────────────────────────────────────────────
\usepackage{booktabs}
\usepackage{tabularx}
\usepackage{array}

% ── Graphics & Diagrams ──────────────────────────────────────
\usepackage{tikz}
\usetikzlibrary{arrows.meta, positioning, shapes.geometric}
\usepackage{graphicx}
\usepackage{float}

% ── Colors ───────────────────────────────────────────────────
\usepackage{xcolor}
\definecolor{vertexcolor}{RGB}{50,  100, 200}
\definecolor{centercolor}{RGB}{200, 80,  50}
\definecolor{yellowpiece}{RGB}{240, 200, 0}
\definecolor{greenpiece} {RGB}{50,  180, 80}

% ── Author / Affiliation (chuẩn cho multi-author, dùng được cả 1 tác giả) ──
\usepackage{authblk}

% ── Hyperlinks ───────────────────────────────────────────────
\usepackage{hyperref}
\hypersetup{
  colorlinks = true,
  linkcolor  = blue,
  urlcolor   = blue,
  citecolor  = blue,
  pdftitle   = {Tactical Formation: Tro Choi Co Chien Thuat Truou Tuong
                Danh Cho 2-4 Nguoi Choi},
  pdfauthor  = {[TEN TAC GIA]}
}

% ── Code Listings ────────────────────────────────────────────
\usepackage{listings}
\lstset{basicstyle=\ttfamily\small, breaklines=true,
        frame=single, captionpos=b}

% ── Misc ─────────────────────────────────────────────────────
\usepackage{enumitem}
\usepackage{caption}
\usepackage{subcaption}
\usepackage{setspace}
\onehalfspacing

% ─────────────────────────────────────────────────────────────
%  PHẦN 1 — TITLE
%  Chuẩn paper khoa học gồm: tiêu đề · tác giả · đơn vị ·
%  email · ngày · abstract · keywords.
%  Tất cả xuất hiện TRƯỚC section đầu tiên.
% ─────────────────────────────────────────────────────────────

% 1a. Tiêu đề
%     Quy tắc: rõ ràng, phản ánh đúng nội dung, không quá 20 từ.
%     Không nhồi 2 ngôn ngữ vào cùng \title{}.
%     Dùng dấu hai chấm để phân tách tên game và mô tả.
\title{%
  \textbf{Tactical Formation:}\\[0.3em]
  \textbf{Trò Chơi Cờ Chiến Thuật Trừu Tượng Dành Cho 2--4 Người Chơi}\\[0.3em]
  \textbf{Với Cơ Chế Di Chuyển Và Lật Quân Kết Hợp}
}

% 1b. Tác giả
%     Chuẩn: tên thật, đơn vị (Independent Researcher nếu độc lập),
%     email liên hệ dùng \thanks{} hoặc footnote.
%     authblk cho phép thêm nhiều tác giả sau nếu cần.
\author{%
  [TÊN TÁC GIẢ]%
  \thanks{Tác giả liên hệ: \href{mailto:quyetops@gmail.com}{quyetops@gmail.com}}
}
\affil{Independent Researcher, Việt Nam}

% 1c. Ngày
%     Chuẩn: chỉ ghi tháng và năm — không ghi "bản nháp".
%     Ghi chú nội bộ để trong comment, không xuất ra paper.
%     TODO: Đổi thành ngày cụ thể khi nộp, ví dụ: \date{Tháng 5, 2025}
\date{\today}

% ─────────────────────────────────────────────────────────────
\begin{document}
% ─────────────────────────────────────────────────────────────

\maketitle

% ─────────────────────────────────────────────────────────────
% 1d. Abstract (Tóm tắt)
%     Chuẩn: 150–250 từ · MỘT đoạn liền · không dùng bullet.
%     Phải trả lời đủ 4 câu hỏi:
%       (1) Vấn đề / bối cảnh là gì?
%       (2) Tôi đã làm gì?
%       (3) Kết quả / đóng góp chính là gì?
%       (4) Ý nghĩa / ứng dụng là gì?
% ─────────────────────────────────────────────────────────────
\begin{abstract}
% (1) Vấn đề / bối cảnh
Các trò chơi cờ trừu tượng kinh điển như Cờ Đam và Reversi mỗi loại chỉ
khai thác một cơ chế tương tác duy nhất và đều giới hạn ở hai người chơi,
tạo ra khoảng trống thiết kế chưa được khám phá khi kết hợp hai cơ chế
đó trong môi trường nhiều người.
% (2) Tôi đã làm gì
Bài báo này giới thiệu \textbf{Tactical Formation}, một trò chơi cờ chiến
thuật trừu tượng mới kết hợp cơ chế \emph{di chuyển bằng cách nhảy qua quân}
từ Cờ Đam với cơ chế \emph{bắt quân bằng kẹp sandwich} từ Reversi,
và mở rộng sang 2, 3 và 4 người chơi đồng thời.
Điểm cốt lõi là quân bị nhảy qua \emph{không bị loại} khỏi bàn mà vẫn tồn tại,
tạo ra căng thẳng chiến thuật chưa xuất hiện trong cả hai game nguồn.
Bàn cờ được thiết kế với hai lớp vị trí đồng thời---đỉnh (vertex) và tâm
(center)---mỗi loại có tập hướng di chuyển riêng, làm tăng đáng kể chiều
sâu chiến thuật trên cùng một mặt bàn.
% (3) Kết quả / đóng góp
Chúng tôi trình bày đặc tả đầy đủ luật chơi cho cả ba biến thể,
và xây dựng AI đồng hành sử dụng Paranoid Minimax với cắt tỉa Alpha-Beta
thay cho Monte Carlo Tree Search---vốn không hiệu quả trong lớp game
Reversi-like do tín hiệu rollout quá nhiễu.
% (4) Ý nghĩa
Phần mềm được phát hành mã nguồn mở, cung cấp nền tảng để cộng đồng
chơi, nghiên cứu, và phát triển các biến thể tiếp theo.
\end{abstract}

% ─────────────────────────────────────────────────────────────
% 1e. Keywords (Từ khóa)
%     Chuẩn: 4–6 từ khóa · viết thường · cách nhau bằng dấu chấm phẩy.
%     Đặt NGAY SAU \end{abstract} — bắt buộc phải có trên arXiv.
% ─────────────────────────────────────────────────────────────
\vspace{0.4em}
\noindent\textbf{Từ khóa:}
trò chơi cờ trừu tượng;
thiết kế game nhiều người;
Paranoid Minimax;
biến thể Reversi;
Alpha-Beta pruning;
cờ chiến thuật.

\vspace{0.8em}
\hrule
\vspace{0.4em}

\tableofcontents
\newpage

% ─────────────────────────────────────────────────────────────
%  PHẦN 2 — INTRODUCTION (Giới Thiệu)
% ─────────────────────────────────────────────────────────────
\section{Giới Thiệu}

Các trò chơi cờ trừu tượng đã tồn tại hàng nghìn năm, từ Cờ Vây của
Trung Quốc đến Cờ Đam châu Âu và Reversi ra đời cuối thế kỷ~19.
Mỗi game định nghĩa một không gian chiến thuật riêng biệt qua tập luật
tối giản của nó.

\textbf{Tactical Formation} ra đời từ câu hỏi thiết kế: \emph{điều gì xảy
ra nếu ta kết hợp cơ chế di chuyển của Cờ Đam với cơ chế bắt quân của
Reversi, nhưng không loại bỏ quân bị nhảy qua, và mở rộng sang nhiều
người chơi?}

Sự kết hợp này tạo ra một căng thẳng chiến thuật mới. Trong Cờ Đam,
người chơi nhảy qua quân để \emph{ăn} nó---quân biến mất. Trong Reversi,
người chơi đặt quân mới để \emph{kẹp} quân đối thủ---quân bị lật màu.
Trong Tactical Formation, người chơi nhảy qua một quân bất kỳ (kể cả quân
của chính mình) để di chuyển, rồi kẹp sandwich từ vị trí mới---quân bị
nhảy qua vẫn còn trên bàn và có thể trở thành vũ khí hoặc rào cản ở lượt
tiếp theo.

Kết quả là mỗi nước đi mang hai chiều tác dụng độc lập: vị trí mới của
quân (ảnh hưởng đến di chuyển tương lai) và sandwich tạo ra từ vị trí đó
(ảnh hưởng đến lật quân ngay lập tức). Người chơi phải cân nhắc cả hai
chiều đồng thời, tạo ra chiều sâu chiến thuật khác biệt so với cả hai
nguồn cảm hứng.

Bài báo có các đóng góp sau:
\begin{enumerate}[label=(\roman*)]
  \item Đặc tả đầy đủ và chính xác luật chơi Tactical Formation cho
        2--4 người chơi.
  \item Mô tả cấu trúc bàn cờ hai lớp (đỉnh và tâm) và phân tích ảnh
        hưởng của nó đến không gian chiến thuật.
  \item Xây dựng AI đồng hành dựa trên Paranoid Minimax, kèm phân tích
        tại sao MCTS không phù hợp cho lớp game này.
  \item Phát hành phần mềm mã nguồn mở để cộng đồng chơi, nghiên cứu
        và phát triển tiếp.
\end{enumerate}

Phần còn lại của bài báo được tổ chức như sau. Phần~\ref{sec:related}
khảo sát các game liên quan. Phần~\ref{sec:rules} đặc tả đầy đủ luật
chơi. Phần~\ref{sec:variants} mô tả các biến thể nhiều người.
Phần~\ref{sec:ai} trình bày thuật toán AI đồng hành.
Phần~\ref{sec:impl} mô tả phần mềm. Phần~\ref{sec:conclusion}
kết luận và đề xuất hướng phát triển tiếp theo.

% ─────────────────────────────────────────────────────────────
%  PHẦN 3 — RELATED WORK (Công Trình Liên Quan)
% ─────────────────────────────────────────────────────────────
\section{Các Game Liên Quan}
\label{sec:related}

\subsection{Cờ Đam (Checkers / Draughts)}

Cờ Đam~\cite{schaeffer2007checkers} là game chiến thuật 2 người trên bàn
8{\texttimes}8. Quân di chuyển theo đường chéo và nhảy qua quân đối thủ
để ăn---quân bị ăn bị loại khỏi bàn. Tactical Formation giữ lại cơ chế
nhảy nhưng bỏ việc loại quân: quân bị nhảy qua trở thành thành phần
chiến lược thay vì bị xóa.

\subsection{Reversi / Othello}

Reversi~\cite{buro1997good} là game 2 người trên bàn 8{\texttimes}8.
Người chơi đặt quân mới và lật tất cả quân đối thủ bị kẹp theo đường
thẳng. Tactical Formation giữ cơ chế lật quân sandwich nhưng thay vì
đặt quân mới, người chơi phải \emph{di chuyển quân hiện có}. Ràng buộc
này giới hạn đáng kể số nước khả dụng và thêm chiều sâu về positioning.

\subsection{So Sánh Tổng Quan}

\begin{table}[H]
  \centering
  \caption{So sánh Tactical Formation với các game cờ liên quan.}
  \label{tab:comparison}
  \begin{tabular}{lccc}
    \toprule
    \textbf{Đặc điểm} & \textbf{Cờ Đam} & \textbf{Reversi}
                      & \textbf{Tactical Formation} \\
    \midrule
    Cách di chuyển   & Nhảy qua quân  & Đặt quân mới & Nhảy qua quân \\
    Quân bị nhảy qua & Bị loại        & ---          & Còn trên bàn  \\
    Cách bắt quân    & Nhảy trực tiếp & Sandwich     & Sandwich sau nhảy \\
    Số người chơi    & 2              & 2            & 2, 3, hoặc 4  \\
    Số lớp vị trí    & 1              & 1            & 2 (đỉnh + tâm) \\
    \bottomrule
  \end{tabular}
\end{table}

% ─────────────────────────────────────────────────────────────
%  PHẦN 4 (PHẦN CHÍNH) — GAME RULES (Luật Chơi)
% ─────────────────────────────────────────────────────────────
\section{Luật Chơi}
\label{sec:rules}

\subsection{Cấu Trúc Bàn Cờ}

Bàn cờ Tactical Formation có \textbf{hai lớp vị trí đồng thời}:

\paragraph{Đỉnh (Vertex).}
Các ô tại tọa độ nguyên $(x,\,y)$. Quân đứng ở đỉnh di chuyển theo
\textbf{8 hướng}: 4 hướng thẳng (bắc, nam, đông, tây) và 4 hướng chéo.

\paragraph{Tâm (Center).}
Các ô tại tọa độ nửa nguyên $(x{+}0.5,\;y{+}0.5)$.
Quân đứng ở tâm chỉ di chuyển theo \textbf{4 hướng chéo}.

Hai lớp hoạt động hoàn toàn độc lập: quân đỉnh chỉ nhảy đến đỉnh khác,
quân tâm chỉ nhảy đến tâm khác.

\begin{figure}[H]
  \centering
  \begin{tikzpicture}[scale=1.1]
    \foreach \x in {0,1,2,3}{
      \foreach \y in {0,1,2,3}{
        \filldraw[fill=vertexcolor!80, draw=vertexcolor!60!black]
          (\x,\y) circle (5pt);
        \ifnum\x<3 \draw[gray!50,thin](\x,\y)--(\x+1,\y); \fi
        \ifnum\y<3 \draw[gray!50,thin](\x,\y)--(\x,\y+1); \fi
      }
    }
    \foreach \x in {0,1,2}{
      \foreach \y in {0,1,2}{
        \filldraw[fill=centercolor!80, draw=centercolor!60!black]
          (\x+0.5,\y+0.5) circle (4pt);
      }
    }
    \node[right, font=\small, color=vertexcolor] at (3.2,1.5)
      {Đỉnh -- Vertex};
    \node[right, font=\small, color=centercolor] at (3.2,1.0)
      {Tâm -- Center};
    \node[above,font=\tiny] at (0,3) {$(0,3)$};
    \node[above,font=\tiny] at (1,3) {$(1,3)$};
    \node[right,font=\tiny,color=centercolor] at (0.55,2.5)
      {$(0.5,2.5)$};
  \end{tikzpicture}
  \caption{Hai lớp vị trí trên bàn cờ 4{\texttimes}4:
           đỉnh (xanh dương) và tâm (đỏ cam).}
  \label{fig:board_layers}
\end{figure}

\begin{table}[H]
  \centering
  \caption{Kích thước bàn cờ và cấu hình theo chế độ chơi.}
  \label{tab:board_size}
  \begin{tabular}{lcccc}
    \toprule
    \textbf{Chế độ} & \textbf{Kích thước} & \textbf{Tọa độ}
    & \textbf{Góc bị loại} & \textbf{Quân/người} \\
    \midrule
    2 người & 5{\texttimes}5 đỉnh
            & $0\text{--}4 \times 0\text{--}4$ & Không & 8 \\
    3 người & 7{\texttimes}7 đỉnh
            & $0\text{--}6 \times 0\text{--}6$ & 4 góc & 12 \\
    4 người & 7{\texttimes}7 đỉnh
            & $0\text{--}6 \times 0\text{--}6$ & 4 góc & 9 \\
    \bottomrule
  \end{tabular}
\end{table}

\subsection{Cách Di Chuyển}

Để di chuyển một quân, cần thỏa mãn đồng thời ba điều kiện:
\begin{enumerate}
  \item Có \textbf{ít nhất một quân kề} (bất kỳ màu nào) ngay cạnh theo
        hướng di chuyển.
  \item Ô đích---phía bên kia quân kề---phải \textbf{trống}.
  \item Ô đích phải cùng loại: đỉnh nhảy đến đỉnh, tâm nhảy đến tâm.
\end{enumerate}
Quân nhảy qua quân kề và đáp xuống ô đích.
\textbf{Quân bị nhảy qua không bị loại khỏi bàn cờ.}

\begin{figure}[H]
  \centering
  \begin{tikzpicture}[scale=1.2]
    \node[font=\small,above] at (1,1.2) {\textbf{Trước}};
    \filldraw[fill=yellowpiece,draw=black](0,0) circle(10pt)
      node[font=\tiny,white]{V};
    \filldraw[fill=greenpiece, draw=black](1,0) circle(10pt)
      node[font=\tiny,white]{X};
    \draw[gray,dashed](2,0) circle(10pt);
    \draw[gray!40](-0.5,0)--(2.5,0);
    \draw[-{Stealth},thick,color=yellowpiece]
      (0.15,0.18) to[bend left=35] (1.85,0.18);
    %
    \begin{scope}[xshift=5cm]
      \node[font=\small,above] at (1,1.2) {\textbf{Sau}};
      \draw[gray,dashed](0,0) circle(10pt);
      \filldraw[fill=greenpiece, draw=black](1,0) circle(10pt)
        node[font=\tiny,white]{X};
      \filldraw[fill=yellowpiece,draw=black](2,0) circle(10pt)
        node[font=\tiny,white]{V};
      \draw[gray!40](-0.5,0)--(2.5,0);
    \end{scope}
  \end{tikzpicture}
  \caption{Quân Vàng nhảy qua quân Xanh và đến ô bên phải.
           Quân Xanh \emph{vẫn còn} trên bàn cờ.}
  \label{fig:movement}
\end{figure}

\subsection{Cách Bắt Quân --- Lật Màu}

Ngay sau khi quân di chuyển đến ô mới, kiểm tra sandwich theo tất cả
các hướng từ ô đó. Một hướng tạo sandwich khi:
\begin{itemize}
  \item Có một hoặc nhiều quân đối thủ \emph{liên tiếp không gián đoạn}
        theo hướng đó, \textbf{và}
  \item Ngay cuối chuỗi đó là một quân \emph{cùng màu} với người vừa
        di chuyển.
\end{itemize}
Tất cả quân đối thủ trong chuỗi đó \textbf{đổi màu} thành màu của người
di chuyển. Nếu nhiều hướng đều tạo sandwich đồng thời, chỉ hướng có
\textbf{số quân bị lật nhiều nhất} được thực hiện. Nếu không có sandwich,
nước đi vẫn hợp lệ.

\begin{figure}[H]
  \centering
  \begin{tikzpicture}[scale=1.1]
    \node[font=\small,above] at (2,0.9) {\textbf{Sandwich hàng ngang}};
    \filldraw[fill=yellowpiece,draw=black](0,0) circle(9pt)
      node[font=\tiny,white]{V};
    \filldraw[fill=greenpiece, draw=black](1,0) circle(9pt)
      node[font=\tiny,white]{X};
    \filldraw[fill=greenpiece, draw=black](2,0) circle(9pt)
      node[font=\tiny,white]{X};
    \filldraw[fill=greenpiece, draw=black](3,0) circle(9pt)
      node[font=\tiny,white]{X};
    \filldraw[fill=yellowpiece,draw=black](4,0) circle(9pt)
      node[font=\tiny,white]{V};
    \draw[gray!40](-0.5,0)--(4.5,0);
    \draw[{Stealth}-{Stealth},red,thick](0.8,-0.5)--(3.2,-0.5);
    \node[below,font=\small,red] at (2,-0.55) {bị lật thành Vàng};
    \node[above,font=\tiny,color=yellowpiece] at (4,0.12) {vừa nhảy};
  \end{tikzpicture}
  \caption{Ba quân Xanh bị kẹp giữa hai quân Vàng và đổi sang màu Vàng.}
  \label{fig:sandwich}
\end{figure}

\subsection{Điều Kiện Thắng Thua}

Một người chơi \textbf{bị loại} khi rơi vào một trong hai trạng thái:
\begin{enumerate}
  \item \textbf{Hết quân}: tất cả quân của họ đã bị lật sang màu khác.
  \item \textbf{Hết nước đi}: không còn quân nào có thể di chuyển hợp lệ.
\end{enumerate}
\textbf{Người cuối cùng còn lại thắng.}

Để tránh hòa kéo dài, nếu sau 150 lượt (3 người) hoặc 200 lượt (4 người)
không có ai bị loại, ván đấu kết thúc và người có \textbf{nhiều quân nhất}
thắng.

% ─────────────────────────────────────────────────────────────
%  PHẦN 5 — VARIANTS (Các Biến Thể Nhiều Người)
% ─────────────────────────────────────────────────────────────
\section{Các Biến Thể Nhiều Người}
\label{sec:variants}

\subsection{Chế Độ 2 Người}
Bàn cờ 5{\texttimes}5 đỉnh. Mỗi bên 8 quân (4 đỉnh + 4 tâm).
Màu: Vàng và Xanh lá.

\subsection{Chế Độ 3 Người}
Bàn cờ 7{\texttimes}7 với 4 góc bị loại tạo hình gần bát giác.
Mỗi bên 12 quân. Màu: Vàng, Xanh lá, Đỏ. Chế độ này giới thiệu
\emph{động lực liên minh}: người chơi phải cân nhắc không chỉ tấn công
đối thủ trực tiếp mà còn tránh tạo lợi thế cho đối thủ thứ ba.

\subsection{Chế Độ 4 Người}
Cùng bàn 7{\texttimes}7 góc cắt, mỗi bên 9 quân.
Màu: Vàng, Xanh lá, Đỏ, Tím. Bốn người chơi tạo ra mạng lưới quan hệ
phức tạp; chiến thắng thường đến từ việc giữ vị trí trung lập rồi ra
đòn quyết định khi các đối thủ đã tiêu hao lẫn nhau.

% ─────────────────────────────────────────────────────────────
%  PHẦN 6 (KỸ THUẬT) — AI COMPANION (AI Đồng Hành)
% ─────────────────────────────────────────────────────────────
\section{AI Đồng Hành}
\label{sec:ai}

\subsection{Lý Do Cần AI}

Một trong những rào cản với game cờ trừu tượng là cần có đối thủ.
AI đồng hành cho phép người chơi luyện tập một mình và cảm nhận chiều
sâu chiến thuật mà không cần chờ người khác.

\subsection{Tại Sao Không Dùng MCTS}

Monte Carlo Tree Search (MCTS)~\cite{coulom2006efficient} thường được
chọn cho game nhiều người vì không cần hàm đánh giá tường minh. Tuy nhiên,
trong các game Reversi-like, MCTS hoạt động kém vì:
\begin{itemize}
  \item Một nước đi có thể lật 5--10 quân lập tức---kết quả một ván ngẫu
        nhiên gần như không liên quan đến chất lượng nước đầu tiên.
  \item Tín hiệu rollout quá nhiễu, UCT không phân biệt được nước tốt
        và nước xấu.
  \item Với budget 500--800 simulations, MCTS chỉ explore được 1--2 lớp
        đầu, tương đương greedy 1-ply.
\end{itemize}
Nghiên cứu trên Othello~\cite{buro1997good} cho thấy Alpha-Beta depth=4
đánh bại MCTS với 100\,000 simulations vì hàm đánh giá chính xác hơn
rollout ngẫu nhiên.

\subsection{Paranoid Minimax}

Với game 2 người, Minimax Alpha-Beta~\cite{knuth1975alphabeta} là đủ.
Với 3--4 người, \textbf{Paranoid Minimax}~\cite{sturtevant2008analysis}
được áp dụng: bot xem \emph{tất cả người chơi còn lại như một liên minh}
chống lại mình, đưa bài toán $N$-người về bài toán max-min 2 chiều và cho
phép dùng cắt tỉa Alpha-Beta đầy đủ.

\paragraph{Hàm đánh giá.}
\begin{equation}
  \mathrm{eval}(\text{pos},\, c)
    = \bigl(\#\text{quân}_c - \#\text{quân}_{\neg c}\bigr) \times 10
      + \#\text{nước đi}_c
  \label{eq:eval}
\end{equation}
Thành phần đầu là ưu thế quân số (trọng số 10);
thành phần sau là tính linh hoạt (mobility).

\paragraph{Thưởng kết thúc nhanh.}
\begin{equation}
  \mathrm{score}_{\text{terminal}} = \pm\,(10000 + d)
\end{equation}
trong đó $d$ là số lớp còn lại trong cây tìm kiếm, khuyến khích bot
\emph{thắng nhanh} thay vì kéo dài ván.

\begin{table}[H]
  \centering
  \caption{Tham số AI theo chế độ chơi.}
  \label{tab:ai_params}
  \begin{tabular}{lccc}
    \toprule
    \textbf{Chế độ} & \textbf{Thuật toán} & \textbf{Độ sâu}
    & \textbf{Thời gian/nước} \\
    \midrule
    2 người & Minimax Alpha-Beta & 6 & $<0.3$\,s \\
    3 người & Paranoid Minimax   & 4 & $<0.5$\,s \\
    4 người & Paranoid Minimax   & 3 & $<0.5$\,s \\
    \bottomrule
  \end{tabular}
\end{table}

AI chạy trong background thread (\texttt{asyncio.run\_in\_executor})
để không làm đứng giao diện trong khi tính toán.

% ─────────────────────────────────────────────────────────────
%  PHẦN 7a — IMPLEMENTATION (Cài Đặt và Phần Mềm)
% ─────────────────────────────────────────────────────────────
\section{Cài Đặt và Phần Mềm}
\label{sec:impl}

Game được viết bằng Python~3 với thư viện Pygame~2. Mã nguồn mở tại:
\url{https://github.com/[username]/Tactical-Formation-The-Board-Game}.

\begin{lstlisting}[language=bash, caption={Cài đặt và chạy game.}]
python -m venv .venv
source .venv/bin/activate   # Linux / macOS
pip install pygame
python main.py
\end{lstlisting}

Phiên bản web chạy qua \texttt{pygbag}~\cite{pygbag}:
\begin{lstlisting}[language=bash, caption={Chạy trên trình duyệt.}]
pip install pygbag
pygbag main.py
# Mở trình duyệt tại http://localhost:8000
\end{lstlisting}

% ─────────────────────────────────────────────────────────────
%  PHẦN 8 — CONCLUSION (Kết Luận)
% ─────────────────────────────────────────────────────────────
\section{Kết Luận}
\label{sec:conclusion}

% 8a. Tóm tắt đóng góp (KHÔNG copy lại abstract)
Bài báo đã trình bày Tactical Formation---một trò chơi cờ chiến thuật
trừu tượng mới kết hợp di chuyển kiểu Cờ Đam và bắt quân kiểu Reversi
trên bàn cờ hai lớp vị trí, hỗ trợ 2--4 người chơi. Luật chơi được
đặc tả đầy đủ và nhất quán cho cả ba biến thể. AI đồng hành dựa trên
Paranoid Minimax hoạt động hiệu quả trong thời gian thực và cho phép
người chơi tự luyện tập khi không có đối thủ.

% 8b. Hạn chế (Limitations) — bắt buộc phải có trong paper chuẩn
\paragraph{Hạn chế hiện tại.}
Phần mềm hiện chưa có chế độ online nhiều người. Hàm đánh giá của AI
chỉ dùng hai yếu tố (quân số và mobility); các yếu tố chiến lược sâu hơn
như kiểm soát trung tâm hoặc đe dọa nhiều hướng chưa được mô hình hóa.
Ngoài ra, chưa có nghiên cứu thực nghiệm định lượng về độ cân bằng giữa
các màu trong chế độ 3 và 4 người.

% 8c. Hướng phát triển tiếp theo (Future Work) — bắt buộc phải có
\paragraph{Hướng phát triển tiếp theo.}
Các hướng nghiên cứu và phát triển tiếp theo bao gồm:
\begin{enumerate}[label=(\roman*)]
  \item Phân tích định lượng độ cân bằng game qua mô phỏng bot-vs-bot
        hàng trăm ván.
  \item Cải thiện hàm đánh giá AI với các đặc trưng chiến lược phức
        tạp hơn, hoặc học hàm đánh giá qua self-play.
  \item Phát triển chế độ online để hỗ trợ nhiều người chơi từ xa.
  \item Khảo sát các biến thể luật chơi, ví dụ cho phép chọn hướng
        sandwich thay vì bắt buộc lấy hướng nhiều nhất.
\end{enumerate}

% ─────────────────────────────────────────────────────────────
%  REFERENCES
% ─────────────────────────────────────────────────────────────
\begin{thebibliography}{9}

\bibitem{schaeffer2007checkers}
  J.~Schaeffer et al.,
  ``Checkers is solved,''
  \textit{Science}, vol.~317, no.~5844, pp.~1518--1522, 2007.

\bibitem{buro1997good}
  M.~Buro,
  ``How machines have learned to play Othello,''
  \textit{IEEE Intelligent Systems}, vol.~14, no.~6, pp.~12--14, 1999.

\bibitem{coulom2006efficient}
  R.~Coulom,
  ``Efficient selectivity and backup operators in Monte-Carlo tree
  search,'' in \textit{Proc. 5th Int. Conf. Computers and Games},
  LNCS vol.~4630, pp.~72--83, 2007.

\bibitem{knuth1975alphabeta}
  D.~E.~Knuth and R.~W.~Moore,
  ``An analysis of alpha-beta pruning,''
  \textit{Artificial Intelligence}, vol.~6, no.~4, pp.~293--326, 1975.

\bibitem{sturtevant2008analysis}
  N.~R.~Sturtevant,
  ``An analysis of UCT in multi-player games,''
  in \textit{Proc. 6th Int. Conf. Computers and Games},
  LNCS vol.~5131, pp.~37--49, 2008.

\bibitem{pygbag}
  P.~Barateau,
  ``pygbag: Python/Pygame WebAssembly packager,''
  \url{https://pygame-web.github.io}, 2023.

\end{thebibliography}

\end{document}
